\documentclass[8pt,landscape]{article}
\usepackage{multicol}
\usepackage{amsmath}
\usepackage{amssymb}
\usepackage{amsthm}
\usepackage{commath}
\usepackage[landscape]{geometry}
\usepackage{enumitem}
\usepackage{graphicx}
\usepackage{mdframed}
\usepackage{algorithm}
\usepackage{algpseudocode}
\usepackage{tikz}
\usepackage{mathrsfs}
\usepackage{titlesec}
\titleformat*{\section}{\small\bfseries}
\titleformat*{\subsection}{\footnotesize\bfseries}

\geometry{top=0.5in,bottom=0.5in,left=0.5in,right=0.5in}

\begin{document}
	\small 
	\begin{multicols*}{4}
		\section*{Module 1: Introduction to DBMS}
		\begin{itemize}[leftmargin=*]
      \item database: collection of data items stored systematically \& permanently in a computer so that programs can consult it to answer questions to get information that can be used to make decisions
      \item benefits of a DBMS: generality \& declarativity, efficiency \& scale, concurrency, consistency \& reliability
      \item relational data model: model data independently of how it will be used, describes data minimally and mathematically
      \item relation: a set of tuples with attributes
      \item schema: sets the structure of the table
      \item arity/degree: number of columns
      \item cardinality: number of rows in current instance
      \item instance of a table (non-header rows) changes quickly over time as the real-world changes
      \item DDL: data definition language, specifies database schema, attributes and their domains, keys (primary or unique), and foreign keys
      \item DML: data manipulation language based on relational algebra and calculus (insert statements are DML)
      \item order in which tuples are inserted into tables matters due to referential integrity (foreign key constraints)
      \item useless query: return the empty result for every instance
      \item common table expressions: creates a temporary table for duration of the query
      \item table expressions: embedded in a query where a table is expected
      \item temporary tables: require permissions, exist for duration of session
		\end{itemize}
    \section*{Module 2: Advanced SQL}
		\begin{itemize}[leftmargin=*]
      \item why sql over excel? (1. performed over larger than memory datasets, 2. data can be combined between tables, 3. optimized to run fast)
      \item set vs bag: a set does not contain duplicates while a bag may contain duplicates
      \item \texttt{SELECT..FROM..WHERE} use bag semantics while \texttt{UNION,INTERSECT,MINUS} use set semantics unless \texttt{ALL} is used
      \item union compatible: two relations have the same number attributes with the same domains
      \item subquery: occurs in \texttt{WHERE} clause, can be correlated (uses variables in outer query) or uncorrelated (independent of outer query)
      \item three state logic: T=1, F=0, U=0.5 where:
        \begin{itemize}
          \item AND: calculate \texttt{MIN(X,Y)}
          \item OR: calculate \texttt{MAX(X,Y)}
          \item NOT: calculate \texttt{1-X}
        \end{itemize}
          \item for comparison with x=NULL we get unknown, arithmetic expressions would get NULL
          \item for aggregating we ignore NULL values except for \texttt{COUNT}
          \item \texttt{SELECT DISTINCT} treats all NULLs as equal
          \item \texttt{COUNT DISTINCT} does not count NULLs
          \item \texttt{GROUP BY} treats NULLs as equal
          \item Set operations treat NULLs as equal unless we use \texttt{ALL}
          \item \texttt{OUTER JOIN}: includes all tuples in the left, right, or both tables even if they do not participate in a join.
          \item cardinality of join result: 0 to product of the sizes of the two tables, if there is a key-foreign relationship, the result is bounded by the size of the table with the foreign key.
          \item view: a named query that can be used in other queries, is computed upon request and is always up-to-date, persists until it is dropped
          \item materialized view: computed once and its results are stored as a table, must be updated as the underlying database changes, important for query optimization
          \item view is updatable if:
            \begin{itemize}
              \item it is defined on a single base table
              \item it uses only selection and projection
              \item there are no aggregates
              \item \texttt{DISTINCT} is not used
              \end{itemize}
		\end{itemize}

    \section*{Module 3: Designing a Database}
		\begin{itemize}[leftmargin=*]
			\end{itemize}
      \section*{Module 4: Relational Design Theory}
		\begin{itemize}[leftmargin=*]
      \item functional dependencies (\(X\rightarrow Y\)): \(X\) and \(Y\) are sets of attributes and are semantic assertions
      \item key: a minimal set of attributes where if the value of all key fields of two tuples are the same, then the tuples are the same
      \item functional dependencies can be enforced by specifying them as a constraint
      \item armstrong's axioms
        \begin{itemize}
          \item reflexivity: if \(Y\subseteq X\) then \(X\rightarrow Y\) (trivial dependencies)
          \item augmentation: if \(X\rightarrow Y\) then \(XW\rightarrow YW\)
          \item transitivity: if \(X\rightarrow Y\) and \(Y\rightarrow Z\) then \(X\rightarrow Z\)
		\end{itemize}
  \item consequences of armstrong's axioms:
    \begin{itemize}
      \item union: if \(X\rightarrow Y\) and \(X\rightarrow Z\) then \(X\rightarrow YZ\)
      \item decomposition: if \(X\rightarrow Y\) and \(Z\subseteq Y\) then \(X\rightarrow Z\)
      \item pseudotransitivity: if \(X\rightarrow Y\) and \(WY\rightarrow Z\) then \(XW\rightarrow Z\)
		\end{itemize}
  \item finding candidate keys: finding subsets \(K\) of \(R\), verifying that \(K\) is a key
    \begin{itemize}
      \item any attribute that is never on the rhs of an fd must be in any primary key
      \item once you have found a candidate key \(K\) you never have to look at supersets of \(K\)
		\end{itemize}
  \item given \(X\) and \(F\) we test if \(X\) is a candidate key for \(R\) by asking if \(X^{+}=R\) and checking that \(X\) is minimal
    \begin{center}
\begin{minipage}{\linewidth}
\begin{algorithmic}[1]
\State $\text{closure} \gets X$
\State Mark all FDs $U \to V$ in $F$ as \textit{unused}
\Repeat
    \If{there is an unused FD $U \to V$ in $F$ such that $U \subseteq \text{closure}$}
        \State Add $V$ to $\text{closure}$ and mark $U \to V$ as \textit{used}
    \EndIf
\Until{no change to $\text{closure}$}
\end{algorithmic}
\end{minipage}
\end{center}
\item 1NF: a relation is in 1NF if all its attributes are atomic (no lists or records)
\item 3NF: A relation is in 3NF if it is in 1NF and for every \(X\rightarrow A\) in F either:
  \begin{itemize}
    \item \(A\in X\) (it is trivial) or
    \item \(X\) is a super key for \(R\) or
    \item \(A\) is a member of some key for \(R\)
    \end{itemize}
  \item BCNF: A relation is in BCNF if for every \(X\rightarrow Y\) in \(F\) either:
    \begin{itemize}
      \item \(Y\) is a subset of \(X\) (i.e. \(X\rightarrow Y\) is trivial) or
      \item \(X\) contains a key for \(R\) (\(X\) is a superkey)
      \end{itemize}
    \item lossless decomposition: \(R_1,\dots,R_{k}\) is a lossless join decomposition of \(R\) wrt an \(FD\) set \(F\) if for every instance \(r\) of \(R\) that satisfies \(F\), projecting \(r\) over each \(R_{i}\) and then joining the resulting instances together gives back the original instance \(r\).
    \item a decomposition is good if it is lossless and it preserves dependencies (\(F=F_1\cup\dots\cup F_k\))
    \item testing for a lossless join: 
      \begin{align*}
        (R_1 \cap R_2)\rightarrow(R_1-R_2)\text{ or}\\
        (R_1 \cap R_2)\rightarrow(R_2-R_1)
      \end{align*}
    \item F is a minimal cover if (intuitively)
      \begin{itemize}
        \item no functional dependency can be removed from \(F\), and
        \item no functional dependency can have its left-hand side simplified by removing an attribute
        \item and still have an equivalent set of dependencies
\end{itemize}
\item 3NF Decomposition algorithm
\begin{center}
\begin{minipage}{\linewidth}
\begin{algorithmic}[1]
\State Let $F$ be a minimal cover, $R$ be a non-3NF schema
\State $i := 0$
\For{each FD $X \to Y$ in $F$}
    \If{none of the schemas $R_j$, $1 \leq j \leq i$, contains $XY$}
        \State $i := i + 1$
        \State $R_i := (X, Y)$
    \EndIf
\EndFor
\If{no schema $R_j$, $1 \leq j \leq i$, contains a candidate key for $R$}
    \State $i := i + 1$
    \State $R_i :=$ any candidate key for $R$
\EndIf
\end{algorithmic}
\end{minipage}
\end{center}
\end{itemize}
    \section*{Module 5: Updates \& Transactions}
		\begin{itemize}[leftmargin=*]
      \item ACID properties of transactions (guaranteed by DBMS):
        \begin{enumerate}
          \item Atomicity: either all of the actions (read/write) of a transaction are executed or none are
          \item Consistency: each transaction executed in isolation keeps the database in a consistent state
          \item Isolation: trnasactions are isolated from the effects of other, currently executing, transactions
          \item Durability: updates persist in the DBMS 
          \end{enumerate}
        \item SQL DML contains \texttt{INSERT, UPDATE, DELETE}
        \item An execution is good it it is serial or equivalent to some serial execution
        \item Undesirable phenomena
          \begin{itemize}
            \item dirty read: a transaction reads data that is written by another \textit{uncommitted} transaction
            \item unrepeatable read: a transaction reads the same data item twice and gets different values
            \item phantom read: a transaction retrieves a set of tuples twice but sees different results, even though it does not modify any of these tuples
            \end{itemize}
          \item SQL isolation levels (strongest to weakest, more expensive to least expensive):
            \begin{itemize}
              \item SERIALIZABLE: complete isolation
              \item REPEATABLE READ: allows phantom reads
              \item READ COMMITTED: allows unrepeatable reads, phantom reads
              \item READ UNCOMMITTED: allows dirty reads, unrepeatable reads, phantom reads (transactions must be read only)
		\end{itemize}
  \item modern systems uses Multi Version Concurrency Control which means that the database maintains multiple versions of the database state
  \item MVCC systems do not implement read uncommitted since avoiding dirty reads comes for free
  \item isolation levels in postgres:
    \begin{enumerate}
      \item read committed (default): Each query in a transaction sees the committed version of the database as of the beginning of the transaction. Changes made by other transactions are visible as soon as they commit.
      \item repeatable read: each query in a transaction sees the committed version of the database as of the beginning of the transaction (so no phantom reads)
      \item serializable: each query in a transaction sees the committed version of the database as of the beginning of the transaction. Read/Write dependencies between transactions are tracked, and if a serialization anomaly was possible
        it aborts one of the transactions (Serializable Snapshot Isolation)
      \item read uncommitted: \textbf{NOT} implemented as dirty reads can \textit{never} occur
      \end{enumerate}
    \item additional undesirable phenomena (not avoided by MVCC and so postgres checks for conflicts when transactions COMMIT)
      \begin{itemize}
        \item Lost Updates: One transaction overwrites the value of data that was written by another transaction
        \item Serialization Anomalies: The execution of a group of transactions is not equivalent to some serial execution of those transactions
  \end{itemize}
\end{itemize}
	\end{multicols*}
\end{document}
